\documentclass{article}
\usepackage[paperwidth=120cm,paperheight=120cm,margin=10mm]{geometry}
\pagestyle{empty}
\usepackage{fontspec}
\usepackage[bidi=basic]{babel}
\babelprovide[import]{english}
\babelprovide[main, import]{persian}
\babelfont{rm}[Renderer=Node, Path=../fonts/, Extension=.ttf, UprightFont=*-Regular, BoldFont=*-Bold]{Vazirmatn}
\babelfont{sf}[Renderer=Node, Path=../fonts/, Extension=.ttf, UprightFont=*-Regular, BoldFont=*-Bold]{Vazirmatn}
\providecommand{\setRTL}{}\providecommand{\setLTR}{}
\providecommand{\RL}[1]{\foreignlanguage{persian}{#1}}
\providecommand{\LR}[1]{\babelsublr{#1}}
\usepackage{tikz}
\usepackage{pgfplots}
\pgfplotsset{compat=1.18}
\usetikzlibrary{arrows.meta, positioning, shapes.geometric, shapes.misc, calc, fit, backgrounds, decorations.pathmorphing, patterns, angles, quotes}
\begin{document}
\setRTL
\babelsublr{\begin{tikzpicture}[>=Stealth, font=\small]

  % ============================================================
  % Slice-selective excitation: RF (apodized sinc) on top,
  % slice-select gradient Gz (main lobe + rephasing lobe) below.
  % Common time axis; effective excitation time at t = Tp/2.
  % ============================================================

  % ---- geometry ----
  \def\x0{0}        % t = 0
  \def\xc{4}        % t = Tp/2 (center of symmetry)
  \def\xp{8}        % t = Tp
  \def\xr{11}       % end of rephasing lobe
  \def\xmid{9.5}    % midpoint of rephasing lobe (=(\xp+\xr)/2)

  % ===================== RF channel (top) =====================
  \begin{scope}[shift={(0,4.2)}]
    % baseline + axis
    \draw[->, thick] (-0.6,0) -- (\xr+0.8,0) node[right] {$t$};
    \draw[->, thick] (\x0,-1.5) -- (\x0,1.7) node[above] {RF: $B_1(t)$};

    % apodized sinc envelope, centered at t = Tp/2 (x=4)
    \draw[very thick, blue!70!black, smooth, samples=200, domain=0:8]
      plot (\x, {1.35*(0.54+0.46*cos(deg(2*pi*(\x-4)/8)))
                 * ( abs(\x-4)<0.01 ? 1 : sin(deg(pi*(\x-4)*0.75))/(deg(pi*(\x-4)*0.75)*pi/180) )});

    % center / effective excitation time
    \draw[dashed, gray] (\xc,-1.5) -- (\xc,1.55);
    \node[gray, anchor=south, font=\footnotesize] at (\xc,1.5) {مرکز تقارن};

    % Tp span
    \draw[<->] (\x0,-1.25) -- (\xp,-1.25);
    \node[fill=white, inner sep=1pt] at (\xc,-1.25) {$T_{\mathrm p}$};
  \end{scope}

  % ===================== Gz channel (bottom) ==================
  \begin{scope}[shift={(0,0)}]
    % axis
    \draw[->, thick] (-0.6,0) -- (\xr+0.8,0) node[right] {$t$};
    \draw[->, thick] (\x0,-1.6) -- (\x0,1.7) node[above] {$G_z(t)$};

    % main slice-select lobe (positive plateau over [0, Tp])
    \def\gA{1.15}  % main amplitude
    \def\gB{1.15}  % rephase amplitude (same |G|, reversed sign)
    \fill[orange!30] (\x0,0) -- (\x0,\gA) -- (\xp,\gA) -- (\xp,0) -- cycle;
    \draw[very thick, orange!80!black]
      (\x0,0) -- (\x0,\gA) -- (\xp,\gA) -- (\xp,0);
    \node[orange!50!black] at (\xc,0.55) {گلبرگ اصلی};

    % rephasing lobe: negative, half the area (duration Tp/2)
    \fill[teal!25] (\xp,0) -- (\xp,-\gB) -- (\xr,-\gB) -- (\xr,0) -- cycle;
    \draw[very thick, teal!70!black]
      (\xp,0) -- (\xp,-\gB) -- (\xr,-\gB) -- (\xr,0);
    \node[teal!50!black, anchor=north, font=\footnotesize] at (\xmid,{-\gB}) {بازفازسازی};

    % span markers
    \draw[<->] (\x0,-1.35) -- (\xp,-1.35);
    \node[fill=white, inner sep=1pt] at (\xc,-1.35) {$T_{\mathrm p}$};
    \draw[<->] (\xp,1.45) -- (\xr,1.45);
    \node[fill=white, inner sep=1pt, font=\footnotesize] at (\xmid,1.45) {$T_{\mathrm p}/2$};

    % alignment line for center
    \draw[dashed, gray] (\xc,-1.6) -- (\xc,\gA+0.1);
  \end{scope}

  % ===================== relation note ========================
  \node[draw, rounded corners, fill=yellow!12, align=center,
        font=\footnotesize] at (\xr+1.4,-2.6)
    {$\displaystyle\int_{\mathrm{rephase}} G_z\,dt
       \approx -\frac{1}{2}\int_0^{T_{\mathrm p}} G_z\,dt$\\[2pt]
     \RL{(مساحت بازفازسازی $\approx$ $-\frac{1}{2}$ مساحت گلبرگ اصلی)}};

  % connector from rephase lobe to note
  \draw[->, gray] (\xmid,-1.25) to[bend right=12] (\xr+0.2,-2.6);

\end{tikzpicture}}
\end{document}
